\SourceLeaf{382}%
\hspace*{\parindent}Je weniger sich hier der anomale Verlauf verkennen lässt, desto interessanter wird die Frage, ob die astronomisch bestimmten Azimute der Dreiecksseiten ähnliche Anomalien zeigen. Nach einem Satz, den \textsc{Laplace} zwar unter besonderen Einschränkungen aufgestellt hat, der sich aber weit verallgemeinern lässt, steht die Konvergenz der Meridiane tatsächlich in einem notwendigen und von der Erdgestalt unabhängigen Zusammenhang mit dem Längenunterschied. sodass sich beim Fortschreiten entlang einer Kette geodätischer Linien die Unregelmäßigkeiten der einen aus denen des anderen \textit{a priori} berechnen lassen. Wie wir berichtet haben, wurden die astronomischen Azimutbestimmungen an den vier Hauptorten Mailand, Turin, Mont Cenis und Colombier mit großer Sorgfalt vorgenommen. Daher haben die Verfasser die Orientierungen an den drei zuletzt genannten Orten mit denen verglichen, die sich ergeben, wenn man die Orientierung in Mailand mithilfe der geodätischen Messungen überträgt. Dabei legten sie die zuvor angegebenen Abmessungen des Erdsphäroids zugrunde. Die Differenzen sind
\begin{center}\begin{tabular}{lr}
für Turin & \(-5''5\)\\
Mont Cenis & \(-51,2\)\\
Colombier & \(-25,2\)
\end{tabular}\end{center}
Auch hier erkennt man also außerordentlich große Anomalien. Berechnet man daraus nach dem erwähnten Satz die Anomalien der Längenunterschiede (indem man durch das Fünfzehnfache des Sinus der Breite von Mailand teilt und das Vorzeichen ändert), so erhält man Werte, die von den unmittelbar gefundenen ganz verschieden sind, nämlich
\begin{center}\begin{tabular}{l|r|r}
& \multicolumn{1}{c|}{Berechnete} & \multicolumn{1}{c}{Unterschied zur}\\
& \multicolumn{1}{c|}{Anomalie} & \multicolumn{1}{c}{beobachteten Anomalie}\\
Turin & \(+0''52\) & \(+2''60\)\\
Mont Cenis & \(+4,81\) & \(+4,11\)\\
Colombier & \(+2,34\) & \(+1,95\)
\end{tabular}\end{center}
Die Verfasser bemerken über \textit{diese} Unterschiede lediglich, dass sie zu groß seien, um sie der Anhäufung der Fehler bei den Winkelmessungen zuschreiben zu können. Sie lassen uns also im Dunkeln darüber, was wir davon halten sollen. Nach unserer Ansicht sind diese drei Zahlen von größter
